\addcontentsline{toc}{section}{СПИСОК ИСПОЛЬЗОВАННЫХ ИСТОЧНИКОВ}
\urlstyle{same}
\renewcommand\UrlFont{\normalfont}
\begin{thebibliography}{9}

    \bibitem{zakon1} Российская Федерация. Законы. О защите прав потребителей: Закон РФ от 7 фервраля 1992г N 2300-1 -- Текст : электронный // КонсультантПлюс: справочно-правовая система. URL: \url{https://www.consultant.ru/document/cons_doc_LAW_305/} (дата обращения: 04.05.2026).
    \bibitem{zakon2} Российская Федерация. Законы. Глава 34 «Аренда». Гражданский кодекс РФ -- Текст : электронный // КонсультантПлюс: справочно-правовая система. — URL: \url{https://www.consultant.ru/document/cons_doc_LAW_9027} (дата обращения: 04.05.2026).
    \bibitem{zakon3} Российская Федерация. Законы. О проведении эксперимента по установлению специального налогового режима "Налог на профессиональный доход": Федеральный закон от 27.11.2018 № 422-ФЗ -- Текст : электронный // КонсультантПлюс: справочно-правовая система. URL: \url{https://www.consultant.ru/document/cons_doc_LAW_311977/} (дата обращения: 04.05.2026).
    \bibitem{zakon4} Российская Федерация. Законы. О персональных данных: Федеральный закон от 27.07.2006 № 152-ФЗ -- Текст : электронный // КонсультантПлюс: справочно-правовая система. URL: \url{https://www.consultant.ru/document/cons_doc_LAW_61801/} (дата обращения: 05.05.2026).
    \bibitem{zakon5} Российская Федерация. Законы. Федеральный закон «Об электронной подписи» от 06.04.2011 № 63-ФЗ // КонсультантПлюс. – URL: \url{https://www.consultant.ru/document/cons_doc_LAW_112701/} (дата обращения 04.05.2026).
    \bibitem{gost100} ГОСТ 19.701-90. Единая система программной документации. Схемы алгоритмов, программ, данных и систем.– Москва : Стандартинформ, 2010. – Текст~: непосредственный.
    \bibitem{gost101} ГОСТ 34.602-2020. Информационные технологии. Комплекс стандартов на автоматизированные системы. Техническое задание на создание автоматизированной системы.– Москва : Стандартинформ, 2020. – Текст~: непосредственный.
    \bibitem{gost102} ГОСТ 34.201-2020. Информационные технологии. Комплекс стандартов на автоматизированные системы. Виды,комплектность и обозначение документов при создании автоматизированных систем.– Москва : Стандарт информ, 2020. – Текст~: непосредственный.
    \bibitem{doc200} Чаплыгин А.А. Основы дипломного проектирования: учебное пособие для студентов направления подготовки 09.03.04 Программная инженерия, направленность (профиль) «Разработка программноинформационных систем» / А.А. Чаплыгин, В.В. Апальков, А.В. Малышев ; Юго-Зап. гос. ун-т. – Курск, 2022. – 102 с. – Библиогр.: 76–79 с. – ISBN 978-5-907407-65-7. – Текст : непосредственный.
    \bibitem{doc202} Котов, А.А. Проблемы разработки систем автоматизированной проверки договоров // Системы анализа и обработки данных. 2025. №1 с. 27-48. – Текст : электронный.
    \bibitem{doc203} Круг, С. Не заставляйте меня думать. Веб-юзабилити и здравый смысл / С. Круг. – Москва : Эксмо, 2017. – 256 с. – ISBN 978-5-699-87233-6. – Текст : непосредственный.
    \bibitem{doc204} Фаулер, М. Рефакторинг. Улучшение существующего кода / М. Фаулер. – Санкт-Петербург : Питер, 2019. – 448 с. – ISBN 978-5-4461-1288-7. – Текст : непосредственный
    \bibitem{doc205} Кальб, И. Объектно-ориентированное программирование с помощью Python / Ирв Кальб ; [перевод с английского М. А. Райтмана]. — Москва : Эксмо, 2024. — 512 с. —- ISBN 978-5-04-186627-3. – Текст : непосредственный.
    \bibitem{doc206} Кочарян, А.С. ДОГОВОР КУПЛИ-ПРОДАЖИ // Вестник НИБ. 2017. №28. – Текст : электронный.
    \bibitem{doc207} Гусак, Е.В  ОПЫТ ВНЕДРЕНИЯ ЭЛЕКТРОННОЙ ПОДПИСИ В КРИТИЧЕСКИХ ДЛЯ БИЗНЕСА СИСТЕМАХ // ТЕНДЕНЦИИ РАЗВИТИЯ НАУКИ И ОБРАЗОВАНИЯ. 2022. №82-2 с. 14-17. – Текст : электронный.
    \bibitem{doc208} Уилсон, Г. Разработка программного обеспечения на примерах с помощью Python / Пер. с англ. О. Перфильев, О. Сайфудинова – М.: Бомбора, 2025 — 402 с. – ISBN 978-5-04-237641-2. – Текст : непосредственный.
    \bibitem{doc209} Меле, А. Django 2 в примерах. Создавайте мощные и надежные веб-приложения Python с нуля / Пер. с англ. Д. В. Плотникова – М.: ДМК Пресс, 2019 — 408 с. – ISBN 978-5-97060-746-6. – Текст : непосредственный.
    \bibitem{doc210} Постолит, А. Python, PyCharm и Django для начинающих / А. Постолит – М.: БХВ-Петербург, 2021 — 464 с. – ISBN 978-5-97756-779-4. – Текст : непосредственный.
    \bibitem{doc211} Шадура, А. Электронная подпись. Просто о сложном / А. Шадура 2019 — 29 с. – ISBN 978-5-0050-2557-9. – Текст : электронный.
    \bibitem{doc212} Дронов, В. А. Django: практика создания Web-сайтов на Python / В. А. Дронов. — СПб.: БХВ-Петербург, 2016. — 528 с. – ISBN 978-5-97750-421-8. – Текст : непосредственный.
    \bibitem{doc213} Бхаргава, А. Грокаем алгоритмы. Иллюстрированное пособие для программистов и любопытствующих. - СПб.: Питер, 2017. - 288 с. – ISBN 978-5-496-02541-6. – Текст : непосредственный.
    \bibitem{doc214} Милл, И. Docker на практике / пер. с англ. Д. А. Беликов. – М.: ДМК Пресс, 2020. – 516 с. – ISBN 978-5-97060-772-5. – Текст : непосредственный.
    \bibitem{doc215} Аттарди, Д. Web API. Сборник рецептов: пер. с англ. - Астана: АЛИСТ, 2025. 304 с. – ISBN 978-601-12-3681-2. - Текст : непосредственный.
    \bibitem{doc216} Мэтиз, Э. Изучаем Python: программирование игр, визуализация данных, веб-приложения. 3-е изд. — СПб.: Питер, 2020 — 512 с. – ISBN 978-5-4461-1528-0. – Текст : непосредственный
    \bibitem{doc217} Мартин, Р. Чистый код. Создание, анализ и рефакторинг / Р. Мартин. – Санкт-Петербург : Питер, 2019. – 464 с. – ISBN 978-5-4461-0960-9. – Текст : непосредственный.
    \bibitem{doc218} Левашов, П. Python с нуля. — СПб.: Питер, 2024. — 448 с. – ISBN 978-5-4461-2145-8. – Текст : непосредственный
    \bibitem{doc219} Константинов, С. С. АРI как искусство: разработка, поддержка, интеграция. - СПб.: БХВ-Петербурr, 2025. - 304 с. – ISBN 978-5-9775-2084-3. – Текст : непосредственный.
    \bibitem{doc220} Прохоренок, Н. А. Python 3. Самое необходимое / Н. А. Прохоренок, В. А. Дронов. — 2-е изд., перераб. и доп. — СПб.: БХВ-Петербург, 2019. — 608 с. – ISBN 978-5-9775-3994-4. – Текст : непосредственный.
    \bibitem{doc221} Арно, Л. Проектирование веб-API / Пер. с англ. Д. А. Беликова.– М.: ДМК Пресс, 2020.– 440 с. – ISBN 978-5-97060-861-6. – Текст : непосредственный.
    \bibitem{doc222} Печурина, Е. В. ДОГОВОР КУПЛИ-ПРОДАЖИ: ИСТОРИЯ И СОВРЕМЕННОСТЬ // Молодой исследователь Дона. 2021. №1 (28). – Текст : электронный.
    \bibitem{doc223} Сирина, О. В. ДОГОВОР АРЕНДЫ В РОССИЙСКОМ ГРАЖДАНСКОМ ПРАВЕ // Universum: экономика и юриспруденция. 2024. №6 (116). – Текст : электронный.
    \bibitem{doc224} Прессман, Р. Разработка программного обеспечения: практическое руководство / Р. Прессман. – Москва : Вильямс, 2016. – 912 с. – ISBN 978-5-8459-2035-5. – Текст : непосредственный
    \bibitem{doc225} Белик, А. Г. Проектирование и архитектура программных систем : учебное пособие / А. Г. Белик, В. Н. Цыганенко. – Омск : ОмГТУ, 2016. – 96 с. –- ISBN 978-5-8149-2258-8. –- Текст : непосредственный.
    \bibitem{doc226} Вигерс, К. Разработка требований к программному обеспечению: [практические приёмы сбора требований и управления ими при разработке программных продуктов : пер. с англ.] / К. Вигерс, Д. Битти. – 3-е изд., доп. – Санкт-Петербург : BHV, 2020. – 736 с. – ISBN 978-5-9775-3348-5. –- Текст : непосредственный.
    \bibitem{doc227} Дронов В. А. Фишерман Л. В. Git. Практическое руководство. Управление и контроль версий в разработке программного обеспечения - Спб.: Наука и техника, 2021. - 304 с. – ISBN 978-5-94387-547-2 – Текст : непосредственный.
    \bibitem{doc228} Юн, Ц. Рецепты Python. Коллекция лучших техник программирования. — СПб.: Питер, 2024. — 512 с. – ISBN 978-5-4461-2156-4. – Текст : непосредственный.
    \bibitem{doc229} Соммервилл, И. Инженерия программного обеспечения / И. Соммервилл. – Москва : Вильямс, 2002. – 624 с. – ISBN 5-8459-0330-0. – Текст : непосредственный.
    \bibitem{doc230} ASGI Documentation: официальный сайт. – URL: \url{https://asgi.readthedocs.io/} (дата обращения 04.05.2026).
    \bibitem{doc231} Redis Documentation: официальный сайт. – URL: \url{https://redis.io/docs/latest/develop/pubsub/} (дата обращения 07.05.2026).
    \bibitem{doc232} Daphne: официальный репозиторий. – URL: \url{https://github.com/django/daphne} (дата обращения 07.05.2026).
    \bibitem{doc233} OpenAPI Specification: официальный сайт. – URL: \url{https://spec.openapis.org/oas/v3.1.0.html} (дата обращения 08.05.2026).
    \bibitem{doc234} drf-yasg: официальная документация. – URL: \url{https://drf-yasg.readthedocs.io/} (дата обращения 09.05.2026).
    \bibitem{doc235} Python Documentation: официальный сайт – URL: \url{https://www.python.org/doc/} (дата обращения: 10.05.2026).
    \bibitem{doc236} Git Documentation: официальный сайт – URL: \url{https://git-scm.com/docs/git} (дата обращения: 11.05.2026).
    \bibitem{doc237} Swagger Documentation: официальный сайт – URL: \url{https://docs.swagger.io/spec.html} (дата обращения: 12.05.2026). 
    \bibitem{doc238} Pycharm Documentation: официальный сайт. – URL:\url{https://www.jetbrains.com/help/pycharm/getting-started.html} (дата обращения: 12.05.2026). 
    \bibitem{doc239} Docker Documentation: официальный сайт. – URL: \url{https://docs.docker.com} (дата обращения: 12.05.2026). 
    \bibitem{doc240} Pillow: официальная документация. – URL: \url{https://pillow.readthedocs.io/} (дата обращения 07.05.2026).
    \bibitem{doc241} Django Documentation: официальный сайт. – URL: \url{https://docs.djangoproject.com/en/6.0/} (дата обращения 15.05.2026).
    \bibitem{doc242} Django REST framework: официальный сайт. – URL: \url{https://www.django-rest-framework.org/} (дата обращения 16.05.2026).
\end{thebibliography}
